\section{The 23-edge case}

Suppose that $m=23$. Let $t$ be the number of degree-two vertices
and keep the notation $R$, $T_{ij}$, and $T$ from the previous section,
now with $\binom{23}{2}=253$. Every degree-two star contains two
edges meeting only at its center. Indeed, if they also met elsewhere,
deleting its center from one of the two edges would leave a five-cover.
Consequently a codegree involving a degree-two vertex is at most one,
and is zero between two distinct degree-two vertices. If $t=11$, every
pair has at least one endpoint with a degree-two vertex, so a codegree
involving degree three is also at most one, by the earlier argument.

\subsection{Profile and star reductions}

Enumerate the six profiles with at least six positive entries, all
between two and nine, and with the sum of the two largest at most 14.
The score must be at least 253 and $t\leq11$. For every pair
of profiles, minimize the collision score across the entire codegree
matrix with given margins and the established cell bounds. After
subtracting the entrywise lower bounds, encode each allowed unit into
a cell by a parallel arc of unit capacity. The successive costs
for adding to a cell of current size $a$ are $a,a+1,\ldots$.
This makes the minimum flow cost plus the initial cost the exact
minimum of $\sum_{u,v}\binom{c_{uv}}2$. Reject if the sum of the
15 minima exceeds $2R$. Repeat with the degree-three cap when $t=11$.

Next consider a particular vertex $v$ of degree $d$. Its codegree
row into another side with profile $q$ is a vector $a$ of sum $d$,
within the same cell bounds. The residual degrees on that side are
$q-a$. The sums of the largest one through four residual
degrees on any side are at most $23-d-9$, $23-d-5$,
$23-d-3$, and $23-d-1$ by the star deletion bounds. The side of $v$
is treated by deleting the entry $d$ from its profile. The residual
cover number is at least five, so at most $23-d$ residual vertices
have degree one. Moreover a row through $v$ uses five other vertices
that cover the residual. At most one of them has residual degree one,
since two such vertices can be replaced by a single vertex covering
their at most two residual edges. Thus the total incidences of the
original star with vertices of residual degree one are at most $d$.
The residual score is at least $\binom{23-d}{2}$, and the cross score
$\sum_{w}c(v,w)(d(w)-c(v,w))$ is at least $d(23-d)$, since every
edge through $v$ meets every edge outside that star.

Enumerate all allowable vectors $a$ independently for each other side.
For each tuple of degree-one count, degree-one incidence in the star,
and residual score, it suffices to retain the largest cross score.
Join these finite tables over the sides. Apply the test to a largest
vertex on every side whose maximum is at least eight. In a second
pass use a threshold of seven and also track collision surplus.
For each candidate $a$ in a pair matrix, remove its fixed row and
minimize the remaining collision score separately by rows and by
columns with the resulting margins and cell bounds. Add the fixed
row cost, take the larger lower bound, and charge any excess above
the unrestricted minimum against $2R$ minus the sum of the pair minima.
Only the five pairs touching the side of $v$ accrue these charges,
so no surplus is counted twice for a test of one vertex.

\subsection{Pair bounds with subset unions and overlap patterns}

Strengthen a pair matrix by applying the bounds on the union of
two through five vertex stars to every choice of vertices from
its two sides. The sum of the selected degrees minus the sum of
entries in the selected rows and columns is the size of the union.
The parallel unit-arc formulation, with these extra linear
constraints, gives a rational relaxation. Every lower bound used
for pruning is validated by an exact rational dual inequality
and rounded up, since the collision count is an integer.

An edge pair meeting in a set $I$ of two, three, or four sides costs
$|I|-1$ from the redundancy budget and contributes one to each required
pair of side indices in $\binom{I}{2}$. A finite recursion chooses
patterns to meet the lower bounds on the 15 side pairs at cost at most
$R$. It branches on a side pair with positive remaining demand and
tries every pattern containing it. When $t=11$, a pattern of size four
must omit at least one side with a degree-two vertex, since at least
one endpoint carries such a vertex and its degree-two side cannot be
among the overlaps.

\begin{center}
\begin{tabular}{l|r}
Necessary conditions applied & Remaining profile systems\\
\hline
Total score and degree-two budget & $1{,}756{,}829$\\
Exact pair transportation & $9{,}614$\\
Degree-three cap when $t=11$ & $5{,}527$\\
First star test & $537$\\
Star test with collision surplus & $126$\\
Subset-union transportation bounds & $72$\\
Overlap pattern budget & $17$
\end{tabular}
\end{center}

\subsection{Row bounds and the last profile}

Every vertex of an edge $e$ has a private neighbor edge, meaning
an edge sharing just that vertex with $e$. Otherwise the vertex could
be omitted from $e$ to obtain a five-cover. If $e$ has a degree-two
vertex, each of its other vertices needs at least two private
neighbors. If such a vertex had at most one, deleting it and the
degree-two vertex would leave at most two uncovered edges, since the
unique neighbor at the degree-two vertex is already private.

For a prospective row of degrees $a_1,\ldots,a_6$, its total excess
from repeated intersections is $\sum_i a_i-28$. Suppose the row has
$b_k$ neighbors with $k$ common vertices. The excess is
$b_2+2b_3+3b_4$. In a side of degree $a_i$, at most $a_i-1$ neighbors
occur there, minus the required number of private neighbors.
Call the resulting capacity for repeated neighbors $u_i$, setting
it to zero at degree two and wherever the degree-three restriction
applies. The $b_4$ neighbors need $4b_4$ incidences in total and
can use at most $\min(u_i,b_4)$ in side $i$. Apply the same inequality
to the combined groups of fourfold and threefold neighbors and to all
repeated neighbors. Enumerating $b_2,b_3,b_4$ gives a lower bound on
$b_2+b_3$ for each allowed degree tuple. Summing over the 23 rows,
with the exact margins for each degree in each side, must give
at most $2(2R-\sum_{i<j}L_{ij})$. The factors of two arise because
each pair has two endpoints. Verified rational linear inequalities
exclude 16 of the 17 profile systems.

The remaining system, in order, is
\[
\begin{gathered}
(8,6,3,2,2,2),\quad (7,7,3,2,2,2),\quad
(8,5,3,3,2,2),\\
(8,5,3,3,2,2),\quad
(8,4,3,3,3,2),\quad (8,3,3,3,3,3).
\end{gathered}
\]
Here $t=11$, $R=23$, and the sum of the strengthened pair minima
is 44. Thus there are seven fourfold pairs and either one threefold
pair or two twofold pairs. Enumeration of the side masks gives no
systems with two twofold pairs and three systems with one threefold pair.

Consider the graph whose vertices are the edges of $H$ and whose
adjacencies are fourfold intersections. If two adjacent edges each
have a degree-two vertex, it must be in the same side. To see this,
remove their four common vertices. An avoiding edge intersects both
in one of two cross combinations on the remaining sides. Both
combinations must occur, since otherwise a fifth vertex suffices
to cover. If the degree-two vertices were on different sides, one
cross combination would contain both, which is impossible.

The fourfold graph has no triangle. A putative triangle has at
least two edges of $H$ with degree-two vertices, in the same side
$i$ by the preceding argument. Any third edge shares neither of
these degree-two labels, since its pair would otherwise meet only
there. All three pairs therefore disagree in side $i$. Their
projections to the other five sides have pairwise Hamming distance
one, so all three projections differ in the same side $j$.
For a triangle edge carrying a degree-two vertex, its unique mate
would have to supply two different $j$-values in the two cross
combinations with the other triangle edges. A single edge cannot do so.

Two neighbors of a row in this triangle-free graph have at least
two agreements, because each differs from the row in two sides.
Their pair therefore must be the unique threefold pair. No row
can have three fourfold neighbors, and at most one row can have
two. If two distinct rows were common to the endpoints of the
threefold pair, those two rows would supply a second pair meeting
in two or three sides. The row with two fourfold neighbors is
outside the unique threefold pair. The two fourfold masks through
that row intersect in precisely the threefold mask, since the two
sets of differing sides must have exactly one common side.

For each of the three overlap-mask systems, enumerate a row's
degree tuple and the repeated-pair labels incident to that row.
Apply the preceding restrictions. For every proper nonempty set
$S$ of side indices, deleting the row's vertices in the other
sides leaves at least $1,3,5,9,14$ edges when $|S|=1,\ldots,5$.
These remaining edges are exactly the row's private neighbors in
$S$ and its repeated neighbors whose entire masks lie in $S$.
This produces $78,71,78$ allowable local rows respectively.

Each mask label needs two endpoints. The degree counts on each
side are fixed. For each of the three systems the verifier checks
an integer linear functional that is nonpositive on every
allowable local row and takes value respectively 55, 53, and 55
on the required aggregate margins. The aggregate therefore cannot
be realized. This excludes 23 edges.
